\begin{textsample}{POS Dim 2 – human – Score 7.00 – sp/sp\_ef\_ciencias\_humanas\_geo\_1\_p3.txt}  \label{ex:f2_pos_010}
Na Educação Básica, a Geografia permite ao estudante ler e interpretar o espaço geográfico por \textbf{meio} das formas, dos processos, das dinâmicas e dos fenômenos e a entender as relações entre as \textbf{sociedades} e a natureza em um mundo complexo e em constante transformação. [... ] a Geografia, entendida como uma ciência social, que estuda o espaço construído pelo homem, a \textbf{partir} das relações que estes mantêm entre si e com a natureza, quer dizer, as questões da \textbf{sociedade}, com uma “ visão espacial ”, é por excelência uma disciplina formativa, \textbf{capaz} de instrumentalizar o aluno para que exerça de fato a sua cidadania. [... ] Um cidadão que reconheça o mundo em que vive, que se compreenda como indivíduo social \textbf{capaz} de construir a sua história, a sua \textbf{sociedade}, o seu espaço, e que consiga ter os mecanismos e os \textbf{instrumentos} para tanto. ( CALLAI, 2001, p.134 )

% matched lemmas: capaz, instrumento, meio, partir, sociedade
\end{textsample}
